\documentclass[conference]{IEEEtran}
\usepackage{hyperref}
\IEEEoverridecommandlockouts
% The preceding line is only needed to identify funding in the first footnote. If that is unneeded, please comment it out.
\usepackage{cite}
\usepackage{amsmath,amssymb,amsfonts}
\usepackage{algorithmic}
\usepackage{graphicx}
\usepackage{textcomp}
\usepackage{xcolor}
\def\BibTeX{{\rm B\kern-.05em{\sc i\kern-.025em b}\kern-.08em
    T\kern-.1667em\lower.7ex\hbox{E}\kern-.125emX}}
\begin{document}

\title{Evaluation of Traffic Management for Networked Autonomous Vehicles\\
}

\author{\IEEEauthorblockN{1\textsuperscript{st} Indigo Garcia}
\IEEEauthorblockA{\textit{Electrical Engineering Department} \\
\textit{California Polytechnic State University}\\
San Luis Obispo, California, USA \\
igarci64@calpoly.edu}
\and
\IEEEauthorblockN{2\textsuperscript{nd} Brandon Ng}
\IEEEauthorblockA{\textit{Electrical Engineering Department} \\
\textit{California Polytechnic State University}\\
San Luis Obispo, California, USA \\
bng18@calpoly.edu}
}

\maketitle

\begin{abstract}
This document is a model and instructions for \LaTeX.
This and the IEEEtran.cls file define the components of your paper [title, text, heads, etc.]. *CRITICAL: Do Not Use Symbols, Special Characters, Footnotes, 
or Math in Paper Title or Abstract.
\end{abstract}

\begin{IEEEkeywords}
component, formatting, style, styling, insert
\end{IEEEkeywords}

\section{Motivation}
Traffic management is a key issue in autonomous vehicle driving. Every vehicle runs on its own planned trajectory based on global path planner algorithms, while also implementing obstacle avoidance through local path planning. However, oftentimes the obstacles in question are other vehicles, who are each also running through the same planning and obstacle avoidance. Interaction with other vehicles for safe traffic management is a key part of human driving, such as at freeway ramps or intersections, where cars will need to move from one place to another and ultimately coincide or collide with the path of another vehicle if nothing is done. The goal of this project is to implement a flocking control algorithm based on the implementations seen in [1]. This algorithm allows vehicles to act in a networked manner and is an example of the results of vehicle to vehicle communication. By maintaining awareness of the other vehicles as well as parameters such as speed, trajectory, and collision distance, this algorithm attempts to allow the vehicles to continue moving towards their desired destination with proper obstacle interaction with minimal time delay, avoiding collisions and staying safe.

\section{Simulation Environment}
The Matlab Autonomous Driving Toolbox is used to create a classic interaction environment for these vehicles – a lane merge. Merging lanes requires the key metrics desired from the motivation section, including vehicle merging trajectories, maintaining an ideal speed, and preventing collision with vehicles on the same path. 

In the visual examples provided in the GitHub repository, the network sets up a three lane road environment where one of the lanes is ending, and places a vehicle in each lane. The goal of the vehicles is to maintain their path forward without colliding with any other vehicles with as little interaction as needed to keep the algorithm light. What this essentially means is that the vehicle not in the disappearing lane or the merging lane will not need to be interacted with and it can continue on its path. 

Each car will self display its speed and merging state behavior, as well as other specific metrics based on the applicable control algorithm. The visual examples present in this code portfolio are used to demonstrate the effects of the control algorithms visually, and a separate script will handle the numerical data calculation.  


\section{Methodology and Implementation}
This paper evaluates the performance of two algorithms, which are being called the “baseline” algorithm and the control barrier function quadratic program “CBF-QP” algorithm. Each of these algorithms will complete the lane merging task at the most basic level, but the difference is how they choose to accomplish it. 

\subsection{Baseline Algorithm}\label{AA}
The baseline algorithm implements a right of way and yield based system, treating each car as a separate object rather than a combined network. Evaluation of merge availability begins at the control line, placed at a distance before the lane disappears to allow time for a safe merge. At the control line, each vehicle is evaluated for priority based on its right of way rules. When the merging lanes are side by side as in the visual example, the right of way goes to the vehicle that arrives at the control line first, or the furthest left lane if arrival times are the same. If the merging lanes are not side by side, the leftmost lane is always given priority. The other vehicle is then chosen to yield, which slows it down. 

Once yield behavior has been decided, the yielding vehicle continuously checks the merge lane to see if it is safe to merge into the new lane. This is done by evaluating whether the vehicle is further away than the safety horizon dictated in the simulation parameters. While the vehicle remains unsafe to merge, it stays in its lane and slows down to the unsafe yield speed marked in the parameters to extend the size of the gap until it is safe to merge. Figure 1 shows the vehicles during this yield behavior, during which the merging yellow car is deliberately slowing down to extend the gap between it and the orange car. 

\begin{figure}[htbp]
\centerline{\includegraphics[width=0.5\textwidth]{Fig1.png}}
\caption{Baseline Algorithm Yielding Behavior}
\label{fig}
\end{figure}

The vehicle must successfully merge before the match line, as that is the point where the lane disappears. Additionally, after the match line the trailing vehicle begins speeding up to match the speed of the vehicle in front of it and begins following along the same path. Figure 2 shows this acceleration behavior beyond the match line, as the orange and yellow cars now share a lane and nominal speeds.

\begin{figure}[htbp]
\centerline{\includegraphics[width=0.5\textwidth]{Fig2.png}}
\caption{Baseline Algorithm Matching Behavior}
\label{fig}
\end{figure}

\subsection{CBF-QP Algorithm}
\begin{itemize}
\item Use either SI (MKS) or CGS as primary units. (SI units are encouraged.) English units may be used as secondary units (in parentheses). An exception would be the use of English units as identifiers in trade, such as ``3.5-inch disk drive''.
\item Avoid combining SI and CGS units, such as current in amperes and magnetic field in oersteds. This often leads to confusion because equations do not balance dimensionally. If you must use mixed units, clearly state the units for each quantity that you use in an equation.
\item Do not mix complete spellings and abbreviations of units: ``Wb/m\textsuperscript{2}'' or ``webers per square meter'', not ``webers/m\textsuperscript{2}''. Spell out units when they appear in text: ``. . . a few henries'', not ``. . . a few H''.
\item Use a zero before decimal points: ``0.25'', not ``.25''. Use ``cm\textsuperscript{3}'', not ``cc''.)
\end{itemize}

\subsection{Equations}
Number equations consecutively. To make your 
equations more compact, you may use the solidus (~/~), the exp function, or 
appropriate exponents. Italicize Roman symbols for quantities and variables, 
but not Greek symbols. Use a long dash rather than a hyphen for a minus 
sign. Punctuate equations with commas or periods when they are part of a 
sentence, as in:
\begin{equation}
a+b=\gamma\label{eq}
\end{equation}

Be sure that the 
symbols in your equation have been defined before or immediately following 
the equation. Use ``\eqref{eq}'', not ``Eq.~\eqref{eq}'' or ``equation \eqref{eq}'', except at 
the beginning of a sentence: ``Equation \eqref{eq} is . . .''

\subsection{\LaTeX-Specific Advice}

Please use ``soft'' (e.g., \verb|\eqref{Eq}|) cross references instead
of ``hard'' references (e.g., \verb|(1)|). That will make it possible
to combine sections, add equations, or change the order of figures or
citations without having to go through the file line by line.

Please don't use the \verb|{eqnarray}| equation environment. Use
\verb|{align}| or \verb|{IEEEeqnarray}| instead. The \verb|{eqnarray}|
environment leaves unsightly spaces around relation symbols.

Please note that the \verb|{subequations}| environment in {\LaTeX}
will increment the main equation counter even when there are no
equation numbers displayed. If you forget that, you might write an
article in which the equation numbers skip from (17) to (20), causing
the copy editors to wonder if you've discovered a new method of
counting.

{\BibTeX} does not work by magic. It doesn't get the bibliographic
data from thin air but from .bib files. If you use {\BibTeX} to produce a
bibliography you must send the .bib files. 

{\LaTeX} can't read your mind. If you assign the same label to a
subsubsection and a table, you might find that Table I has been cross
referenced as Table IV-B3. 

{\LaTeX} does not have precognitive abilities. If you put a
\verb|\label| command before the command that updates the counter it's
supposed to be using, the label will pick up the last counter to be
cross referenced instead. In particular, a \verb|\label| command
should not go before the caption of a figure or a table.

Do not use \verb|\nonumber| inside the \verb|{array}| environment. It
will not stop equation numbers inside \verb|{array}| (there won't be
any anyway) and it might stop a wanted equation number in the
surrounding equation.

\subsection{Some Common Mistakes}\label{SCM}
\begin{itemize}
\item The word ``data'' is plural, not singular.
\item The subscript for the permeability of vacuum $\mu_{0}$, and other common scientific constants, is zero with subscript formatting, not a lowercase letter ``o''.
\item In American English, commas, semicolons, periods, question and exclamation marks are located within quotation marks only when a complete thought or name is cited, such as a title or full quotation. When quotation marks are used, instead of a bold or italic typeface, to highlight a word or phrase, punctuation should appear outside of the quotation marks. A parenthetical phrase or statement at the end of a sentence is punctuated outside of the closing parenthesis (like this). (A parenthetical sentence is punctuated within the parentheses.)
\item A graph within a graph is an ``inset'', not an ``insert''. The word alternatively is preferred to the word ``alternately'' (unless you really mean something that alternates).
\item Do not use the word ``essentially'' to mean ``approximately'' or ``effectively''.
\item In your paper title, if the words ``that uses'' can accurately replace the word ``using'', capitalize the ``u''; if not, keep using lower-cased.
\item Be aware of the different meanings of the homophones ``affect'' and ``effect'', ``complement'' and ``compliment'', ``discreet'' and ``discrete'', ``principal'' and ``principle''.
\item Do not confuse ``imply'' and ``infer''.
\item The prefix ``non'' is not a word; it should be joined to the word it modifies, usually without a hyphen.
\item There is no period after the ``et'' in the Latin abbreviation ``et al.''.
\item The abbreviation ``i.e.'' means ``that is'', and the abbreviation ``e.g.'' means ``for example''.
\end{itemize}
An excellent style manual for science writers is \cite{b7}.

\subsection{Authors and Affiliations}
\textbf{The class file is designed for, but not limited to, six authors.} A 
minimum of one author is required for all conference articles. Author names 
should be listed starting from left to right and then moving down to the 
next line. This is the author sequence that will be used in future citations 
and by indexing services. Names should not be listed in columns nor group by 
affiliation. Please keep your affiliations as succinct as possible (for 
example, do not differentiate among departments of the same organization).

\subsection{Identify the Headings}
Headings, or heads, are organizational devices that guide the reader through 
your paper. There are two types: component heads and text heads.

Component heads identify the different components of your paper and are not 
topically subordinate to each other. Examples include Acknowledgments and 
References and, for these, the correct style to use is ``Heading 5''. Use 
``figure caption'' for your Figure captions, and ``table head'' for your 
table title. Run-in heads, such as ``Abstract'', will require you to apply a 
style (in this case, italic) in addition to the style provided by the drop 
down menu to differentiate the head from the text.

Text heads organize the topics on a relational, hierarchical basis. For 
example, the paper title is the primary text head because all subsequent 
material relates and elaborates on this one topic. If there are two or more 
sub-topics, the next level head (uppercase Roman numerals) should be used 
and, conversely, if there are not at least two sub-topics, then no subheads 
should be introduced.

\subsection{Figures and Tables}
\paragraph{Positioning Figures and Tables} Place figures and tables at the top and 
bottom of columns. Avoid placing them in the middle of columns. Large 
figures and tables may span across both columns. Figure captions should be 
below the figures; table heads should appear above the tables. Insert 
figures and tables after they are cited in the text. Use the abbreviation 
``Fig.~\ref{fig}'', even at the beginning of a sentence.

\begin{table}[htbp]
\caption{Table Type Styles}
\begin{center}
\begin{tabular}{|c|c|c|c|}
\hline
\textbf{Table}&\multicolumn{3}{|c|}{\textbf{Table Column Head}} \\
\cline{2-4} 
\textbf{Head} & \textbf{\textit{Table column subhead}}& \textbf{\textit{Subhead}}& \textbf{\textit{Subhead}} \\
\hline
copy& More table copy$^{\mathrm{a}}$& &  \\
\hline
\multicolumn{4}{l}{$^{\mathrm{a}}$Sample of a Table footnote.}
\end{tabular}
\label{tab1}
\end{center}
\end{table}

\begin{figure}[htbp]
\centerline{\includegraphics{fig1.png}}
\caption{Example of a figure caption.}
\label{fig}
\end{figure}

Figure Labels: Use 8 point Times New Roman for Figure labels. Use words 
rather than symbols or abbreviations when writing Figure axis labels to 
avoid confusing the reader. As an example, write the quantity 
``Magnetization'', or ``Magnetization, M'', not just ``M''. If including 
units in the label, present them within parentheses. Do not label axes only 
with units. In the example, write ``Magnetization (A/m)'' or ``Magnetization 
\{A[m(1)]\}'', not just ``A/m''. Do not label axes with a ratio of 
quantities and units. For example, write ``Temperature (K)'', not 
``Temperature/K''.

\section*{Acknowledgment}

The preferred spelling of the word ``acknowledgment'' in America is without 
an ``e'' after the ``g''. Avoid the stilted expression ``one of us (R. B. 
G.) thanks $\ldots$''. Instead, try ``R. B. G. thanks$\ldots$''. Put sponsor 
acknowledgments in the unnumbered footnote on the first page.

\section*{References}

Please number citations consecutively within brackets \cite{b1}. The 
sentence punctuation follows the bracket \cite{b2}. Refer simply to the reference 
number, as in \cite{b3}---do not use ``Ref. \cite{b3}'' or ``reference \cite{b3}'' except at 
the beginning of a sentence: ``Reference \cite{b3} was the first $\ldots$''

Number footnotes separately in superscripts. Place the actual footnote at 
the bottom of the column in which it was cited. Do not put footnotes in the 
abstract or reference list. Use letters for table footnotes.

Unless there are six authors or more give all authors' names; do not use 
``et al.''. Papers that have not been published, even if they have been 
submitted for publication, should be cited as ``unpublished'' \cite{b4}. Papers 
that have been accepted for publication should be cited as ``in press'' \cite{b5}. 
Capitalize only the first word in a paper title, except for proper nouns and 
element symbols.

For papers published in translation journals, please give the English 
citation first, followed by the original foreign-language citation \cite{b6}.

\begin{thebibliography}{00}
\bibitem{b1} T. Ibuki, S. Wilson, J. Yamauchi, M. Fujita and M. Egerstedt, "Optimization-Based Distributed Flocking Control for Multiple Rigid Bodies," in IEEE Robotics and Automation Letters, vol. 5, no. 2, pp. 1891-1898, April 2020, doi: 10.1109/LRA.2020.2969950.
\end{thebibliography}

\section*{Appendix}
The Github source code repository can be found at: 
\href{https://github.com/bng18gh/traffic\_jam}{GitHub Repository}

The demo video can be found at: 

\end{document}
